\subsection{研究設計}

本研究の最終成果は、精密な数値予測ではなく、段階ベースのギャップ評価である。
本研究の方法論上の価値は、図~\ref{fig:method-workflow}に示すように、幅・深さを回路資源の中間証拠として定義し、文献から抽出した証拠を、社会側要求段階と量子側技術段階のギャップへ変換し、最後に実用候補性の条件として読める形へ整理する手続きにある。
分析は次の手順で行う。

\begin{enumerate}
  \item 既存量子VRP/CVRP/EVRP研究の問題設定、定式化、実行環境、技術段階を整理する。
  \item 各文献から幅、深さ、深さの定義、ゲート数、ショット数、実行可能性、ワークフロー情報を抽出する。
  \item 抽出値を、理論定義、小規模シミュレーション、小規模実機を意識したワークフロー、資源推定などの量子側準備度証拠へ分類する。
  \item 社会側資料からEV・物流システムの社会段階を整理し、問題、データ・ワークフロー、制約、資源、技術段階の各軸に分ける。
  \item 社会側準備度と量子側準備度を対応づけ、対応できない部分をギャップとして記録する。
  \item ギャップを、規模、資源、ハードウェア・ワークフロー、実行可能性、幅・深さのトレードオフとして整理し、実用候補化に必要な証拠条件へ翻訳する。
  \item 文献証拠を順序尺度スコアへ変換し、閾値と判定条件を用いてシナリオ別のSRL--QRL状態遷移を評価する。
\end{enumerate}

この手続きにより、単一の量子ビット数や深さの値をもって実用性を判定するのではなく、どの社会側要求に対して、どの量子側証拠が不足しているかを追跡できる。
したがって図~\ref{fig:method-workflow}は、独立した概念図ではなく、上記手続きの研究内での使い方を示す方法論ワークフローである。

\subsection{指標}

本研究では、量子VRP/CVRP/EVRP研究の実用候補性を直接判定するのではなく、社会側要求に対して量子側証拠がどの段階にあるかを評価する。
そのため、本研究の指標は「性能指標」ではなく、文献から抽出可能な回路資源・ワークフロー証拠として定義する。
主指標は幅と深さであり、ゲート数、ショット数、実行時間、実行可能性、実行環境の情報は、幅・深さの解釈を補助する副指標として用いる。

\begin{table}[t]
\centering
\caption{分析対象文献から抽出する指標}
\label{tab:indicators}
\small
\begin{tabular}{p{0.18\linewidth}p{0.36\linewidth}p{0.34\linewidth}}
\toprule
指標 & 本研究での操作的定義 & 使い方 \\
\midrule
幅 & 問題インスタンスまたは定式化を量子回路へ写像したときに必要となる量子ビット数。QUBO/Isingで1つの二値変数を1量子ビットへ写像する場合は二値変数数と一致する。ただし、最小符号化やHOBOでは一致しない場合がある。 & 図Aの主指標。問題規模が量子回路幅へどう写るかを見る。 \\
深さ & 論文が報告する回路深さまたは層数。QAOA層 $p$、ansatz層 $L$、理論深さ、正規化深さ、トランスパイル後の深さを含む。 & 図Bの主指標。ただし深さの定義を必ず併記し、同一単位の実測値としては扱わない。 \\
深さの定義 & 深さが何を意味するかを示す分類。例として、QAOA層、AOAの位相・ミキサー対、ハードウェア効率型ansatz層、理論上界、トランスパイル後または正規化後の深さを含む。 & 深さの比較可能性を制限する注記。 \\
ゲート・体積指標 & 論文が報告するゲート数、CNOTまたは2量子ビットゲート数、回路体積、量子体積、漸近ゲート複雑性。 & ハードウェア・ワークフローギャップや資源ギャップの補助証拠。 \\
ショット数・実行回数 & ショット数、測定回数、最適化器の実行回数、初期点数、実行時間など。 & ハイブリッドワークフロー負担、実機実行負担、再現性の補助証拠。 \\
実行環境 & シミュレータ、状態ベクトル、QASMシミュレータ、IBM/Rigetti/IonQ実機、資源推定など。 & QRL/ARL型の技術段階分類。 \\
実行可能性・ワークフロー上の警告 & 実行不可能解のサンプリング、実行可能経路の確率、実行時間制約、ノイズ感度、後処理の必要性など。 & 実行可能性ギャップ、ハードウェア・ワークフローギャップの根拠。 \\
\bottomrule
\end{tabular}
\normalsize
\end{table}

\subsection{回路資源の抽出手順}

抽出は、文献単位ではなく、同一文献内のインスタンス、符号化、実行条件ごとに行った。
同じ論文でも、完全符号化と最小符号化、QUBOとHOBO、4ノードと5ノードのように回路資源が異なる場合は別行として記録した。
抽出結果の単位は \texttt{circuit\_resources.csv} の1行であり、付録~\ref{app:extracted-circuit-resources}に本文確認用の表として再掲する。

具体的な抽出手順は次の通りである。

\begin{enumerate}
  \item 対象文献から、問題種類、インスタンス規模、符号化・定式化、実行環境を特定する。
  \item 二値変数数、経路数、顧客数、車両数、レジスタ構造など、幅に対応する記述を探す。
  \item 論文が量子ビット数を直接報告している場合はその値を記録し、式のみの場合は式をそのまま記録する。式から導出した場合は注記に導出根拠を残す。
  \item 深さについては、数値や式だけでなく、それがQAOA層 $p$、ansatz層 $L$、理論深さ、正規化深さ、トランスパイル後の深さのどれかを必ず記録する。
  \item ゲート数、ショット数、実行時間、最適化器の実行回数、実行環境、実行可能性に関する警告が報告されている場合は副指標として記録する。
  \item 図表番号、節番号、または抽出テキスト上の検索手がかりを出典位置として記録する。
  \item 定義、方法論、抽出済み、部分抽出のいずれかの抽出状態を付け、QAOA定義・ベンチマーク方法論・個別応用インスタンスを混同しないようにする。
\end{enumerate}

本研究で抽出したデータは、合計27行である。
このうち、QAOA定義が1行、ベンチマーク方法論が1行、VRPTWが5行、HVRPが3行、CVRP資源推定が2行、TSP符号化理論が4行、VRP-QAOAシミュレーションが3行、ゲート型VRPの実機志向研究が3行、実行可能性を保つCVRP AOAが2行、CVRP/TSP分解研究が3行である。
したがって、図A--Cは、単一種類の実測データではなく、定義、方法論、小規模シミュレーション、実機を意識したワークフロー、資源推定を、深さの定義と技術段階で分けて読ませるための図である。

\subsection{レイヤー分類}

変数や処理は、次のレイヤーに分類する。

\begin{itemize}
  \item 量子実行レイヤー: 幅、深さ、ゲート数、ショット数、エラー率、実行可能なビット列のサンプリング
  \item 符号化・二値定式化レイヤー: QUBO、Ising、HOBO、ペナルティ項、ミキサー、制約を保つ符号化
  \item ハイブリッドワークフローレイヤー: 最適化器の反復、パラメータ更新、トランスパイル、サンプリング、後処理
  \item 古典前処理・最適化レイヤー: クラスタリング、経路生成、分解、古典ベースラインソルバー、実行可能性の修復
  \item 技術社会的な外生レイヤー: EV保有台数、物流需要、充電インフラ、政策、都市構造
  \item 間接的な運用成果レイヤー: 経路品質、稼働率、充電待ち、配送遅延、再最適化頻度
\end{itemize}

\subsection{実用候補性へ向けた変換手順}

図~\ref{fig:method-workflow}に示す変換手続きは、次のように定義する。

\begin{enumerate}
  \item 回路資源の抽出: 分析対象文献から、幅、深さ、ゲート数、ショット数、実行時間、実行可能性、資源推定を抽出する。
  \item 証拠の型分け: 抽出値を、定義値、シミュレーション値、実機を意識したワークフロー値、資源推定値に分ける。
  \item 準備度の対応づけ: 社会側SRLと量子側QRLを対応づけ、どの準備度軸で証拠が不足するかを記録する。
  \item ギャップのラベル付け: 不足部分を、規模ギャップ、資源ギャップ、ハードウェア・ワークフローギャップ、実行可能性ギャップ、幅・深さのトレードオフとして分類する。
  \item 実用上の解釈: ギャップのラベルを、実用候補として次に必要な証拠条件へ翻訳する。
\end{enumerate}

この手続きは、量子優位を直接判定するものではない。
むしろ、既存研究が「実用候補として評価される前に、何を追加で示す必要があるか」を特定するための方法である。

\subsection{準備度の状態遷移評価}

最後に、表Dで得た現状ギャップに対して、技術社会改善シナリオを適用する。
この評価は、量子回路の実行時間、配送性能、費用対効果を予測するものではない。
文献から得られた証拠を0--3の順序尺度スコアとして符号化し、SRL/QRLの段階対応がどの条件で改善するかを調べる状態遷移評価である。

量子側証拠は、資源、実行可能性、ワークフローの3種類に分ける。
資源スコアは幅、深さ、ゲート数、資源推定など、実行可能性スコアは実行可能経路や制約充足に関する証拠、ワークフロースコアはショット数、実行時間、最適化ループ、実行環境の負担などに基づく。
社会側準備度は、問題、データ・ワークフロー、制約、資源、技術段階の5軸を0--3で符号化し、その平均を補助的な社会側準備度スコアとして用いる。

遷移判定では、過大主張を避けるため保守的な閾値と判定条件を用いる。
QRL2には少なくとも基本的な資源証拠を要求し、QRL3には資源証拠と最低限のワークフロー証拠を要求する。
QRL4は資源推定段階として扱うため、資源スコアが高いだけでなく、応用問題に即した資源推定証拠を要求する。
また、実用候補には、資源、実行可能性、ワークフロー、社会側準備度がすべて一定水準以上であることを要求する。
準備度フレームワークやベンダーのロードマップは、QRL定義、解釈、将来シナリオの根拠として用いるが、現在のQRLを直接上げる証拠としては扱わない。

シナリオは、技術主導、社会需要主導、制約共進化、ワークフロー共進化、資源共進化、技術社会の均衡改善の6種類とする。
各シナリオでは、対応するスコアまたは目標SRLを変更し、フェーズS4で固定した閾値・判定条件を用いて、シナリオ別の到達可能QRL、対応状態、残るギャップ、次に必要な証拠を算出する。
したがって、本文の表~\ref{tab:transition-summary}は、性能予測ではなく、どの改善方向であれば実用候補条件に近づくかを示す構造的な評価表である。
全ギャップ・全シナリオの詳細な判定結果は、再現性確認のため付録表~\ref{tab:transition-results}に示す。

\subsection{シナリオメイクの比較設計}

本研究の状態遷移評価は、文献から抽出した証拠、SRL--QRL対応、および明示的なシナリオ仮定に基づく制約付きのシナリオメイクである。
これを、社会側の将来期待やステークホルダーの問題認識から出発する探索的シナリオメイクと比較することで、現在の技術証拠が支持できる範囲と、社会側が期待する応用像とのずれを評価できる。
本研究では、前者を evidence-informed scenario、後者を exploratory scenario と呼ぶ。

Evidence-informed scenario は、抽出済みの幅、深さ、ゲート数、ショット数、実行可能性、ワークフロー証拠、資源推定に加えて、SRL/QRL定義、対象インスタンス規模、社会側制約、技術改善シナリオに関する仮定を明示して構成する。
そのため、このシナリオは文献から直接抽出した事実そのものではなく、抽出証拠と仮定に基づく保守的な現状診断であり、文献証拠が不足する箇所ではQRL不足、判定条件未充足、条件付きと判定される。
一方、exploratory scenario は、既存文献の証拠に直接制約されず、将来のEV・物流社会で期待される機能、運用条件、政策目標、現場側の要求から構成する。
ただし、これは根拠のない予測ではなく、社会側期待にもとづく仮説的シナリオとして扱う。

両者の比較では、次の3種類の差分を確認する。

\begin{enumerate}
  \item 社会側期待が技術証拠を上回る領域: exploratory scenarioでは重要だが、evidence-informed scenarioではQRL不足や判定条件未充足として残る領域。
  \item 技術側証拠はあるが社会側接続が弱い領域: 幅・深さや資源推定は進んでいるが、TMS、配車データ、制約定義、古典ベースラインとの接続が弱い領域。
  \item 両者が一致する有望領域: 社会側期待と量子側証拠が同じ方向を指し、次に必要な実験・実証条件を具体化できる領域。
\end{enumerate}

この比較設計により、本研究のシナリオ分析は単なる将来予測ではなく、文献ベースの制約付き評価と社会側期待の差分を明らかにする方法として位置づけられる。
ワークショップを行う場合には、まず表~\ref{tab:gap-matrix}と表~\ref{tab:transition-summary}に基づくevidence-informed scenarioを提示し、その後、参加者が作成するexploratory scenarioと比較する。
これにより、量子最適化技術に対して社会側が期待する機能と、現在の技術証拠が支持できる範囲とのギャップを明示できる。

本研究の方法論全体は、未来を単一の予測値として推定するものではなく、将来の交通社会に関する複数の可能性を整理し、現在の技術証拠から支持できる範囲を明らかにするforesight型の設計として位置づけられる。
Foresight研究では、入力情報を解釈し、将来像を構成し、それを戦略や行動へ接続するプロセスが重視される\cite{voros2003foresight}。
また、シナリオ研究では、予測的、探索的、規範的なシナリオを区別し、目的に応じて技法を選択する必要がある\cite{borjeson2006scenario,popper2008foresightmethods}。
この観点から、本研究のevidence-informed scenarioは文献証拠と明示的仮定に制約された現状診断型シナリオ、exploratory scenarioは参加者が想定する未来交通社会の枝先を探索するシナリオ、PoC候補状態からの議論は規範的・バックキャスト型シナリオとして扱う。

\subsection{バックキャスト型ワークショップ設計}

表~\ref{tab:gap-matrix}と表~\ref{tab:transition-summary}は、ワークショップで提示するシナリオそのものではなく、evidence-informed scenarioを構成するための現状診断として用いる。
すなわち、文献から抽出した証拠と明示的仮定に基づいて、どのギャップがどの改善シナリオでも残りやすいかを整理する基準線である。
ワークショップでは、この基準線を出発点として、PoC候補状態またはresource estimation候補状態を先に定義し、そこから逆算するバックキャスト型の議論を行う。
この位置づけは、設計したartifactを構築し評価するdesign science researchの考え方と整合する。
本研究におけるartifactは、SRL--QRL gap assessment、evidence-informed baseline、PoC候補状態、resource estimation候補状態を組み合わせたコンセプトデザインである。
ワークショップは、このコンセプトデザインのdemonstrationおよびformative evaluationとして位置づける\cite{hevner2004designscience,peffers2007dsrm,venable2016feds}。

PoC候補状態は、SRLとQRLの単純な合計点だけでは定義しない。
例えば、SRL5とQRL1の組み合わせは合計点が高くても、社会側要求が技術側証拠を大きく上回っている状態であり、PoC候補とは解釈できない。
したがって、本研究では、合計点ではなく最低条件を伴う到達状態としてPoC候補を定義する。
例として、小規模PoC候補は、社会側要求が少なくともSRL3程度であり、量子側証拠がQRL3程度に達し、資源、実行可能性、ワークフロー証拠が一定水準を満たす状態とする。
また、resource estimation候補は、実運用に対応する対象インスタンスが定義され、QRL4相当の資源推定証拠へ接続できる状態とする。

バックキャストでは、参加者に対して次のような問いを提示する。

\begin{enumerate}
  \item このPoC候補状態に到達するために、社会側で先に整えるべきデータ、制約、運用ワークフローは何か。
  \item 量子側で追加すべき証拠は、資源推定、実行可能解率、ショット数・実行時間、古典ベースライン比較のどれか。
  \item 表~\ref{tab:transition-summary}で示されたQRL不足や判定条件未充足を、どの実験・実証・データ整備で解消できるか。
  \item そのPoC候補状態は、参加者が想定する交通社会・EV物流社会にとって意味のあるユースケースか。
\end{enumerate}

この設計により、表~\ref{tab:gap-matrix}と表~\ref{tab:transition-summary}は単独の意思決定表ではなく、バックキャスト型ワークショップで用いるevidence-informed baselineとして位置づけられる。
ワークショップの目的は、これらの表の判定をそのまま受け入れることではなく、PoC候補状態から逆算したときに、社会側・技術側でどの条件を優先的に整備すべきかを検証することである。
加えて、交通社会の未来には複数の枝先があり得るため、ワークショップでは、用意したevidence-informed scenarioを評価するだけでなく、参加者が想定する将来の交通・物流構想を探索する。
その結果は、技術側ではPoC候補、resource estimation対象、feasibility検証、workflow計測、古典ベースライン比較へ接続し、社会側では将来交通社会の構想、必要データ、制約設計、運用workflow要件へ接続する。
このような望ましい将来状態から逆算する議論は、backcastingやparticipatory scenario planningの考え方と整合する\cite{quist2006backcasting}。

\subsection{アンケートによる方法論の検証}

本研究の方法論の意義は、SRL--QRLギャップを分類すること自体ではなく、その分類をPoC、資源推定、実験設計、運用データ整備などの次のエンジニアリング上の行動へ翻訳できる点にある。
したがって、方法論の妥当性は、量子優位や配送性能を直接実証することではなく、提示したギャップ評価が、専門家やステークホルダーにとって次に取るべき行動を明確にするかどうかによって検証できる。

アンケートまたはワークショップでは、参加者にevidence-informed scenarioとexploratory scenarioを提示し、次の観点を評価してもらう。

\begin{enumerate}
  \item ギャップ評価の理解可能性: 表~\ref{tab:gap-matrix}や表~\ref{tab:transition-summary}から、どの社会要求に対してどの量子側証拠が不足しているかを理解できるか。
  \item 次アクションへの接続性: 提示された `next required evidence` や `blocking condition` が、PoC、resource estimation、実機実験、古典ベースライン比較、TMS・配車データ整備などの行動に翻訳できるか。
  \item 優先順位付けへの有用性: どのギャップを先に検証すべきか、どのシナリオをPoC候補とすべきかを判断しやすくなるか。
  \item 社会側期待とのずれの可視化: exploratory scenarioで出た期待やユースケースが、evidence-informed scenarioのどのQRL不足や判定条件未充足に対応するかを確認できるか。
\end{enumerate}

この検証は、方法論が実際の量子計算性能を予測したかを問うものではない。
むしろ、SRL--QRLギャップ評価が、実証前段階における研究計画、PoC設計、資源推定、運用要件定義を支援するかを確認するための検証である。
アンケート項目は、提示したコンセプトデザインの理解可能性、有用性、次アクション性、参加者の業務・研究課題との適合性を測る。
特に、perceived usefulnessやtask-technology fitの考え方を参照し、参加者が本方法論をPoC設計や社会側構想設計に使えると感じたかを測定する\cite{davis1989tam,goodhue1995ttf}。
